\section*{Bab \ref{ch:g10:stats} --- Statistika Deskriptif}

\begin{solution}{exo:g10:stats:1}
Jumlahnya: $0+1+1+2+2+2+3+3+4+4+5+6+7+8+12 = 60$, jadi \omterm{def:g10:stats:mean}{rata-ratanya}
$\frac{60}{15} = 4$. Datanya sudah terurut dan $N = 15$ ganjil:
\omterm{def:g10:stats:median}{mediannya} adalah nilai ke-$8$, yaitu $3$. \omterm{def:g10:stats:quartiles}{Jangkauannya}: $12 - 0 = 12$.
\end{solution}

\begin{solution}{exo:g10:stats:2}
\omterm{def:g10:stats:series}{Frekuensi relatifnya}: $\frac{7}{50} = 0.14$, $\frac{9}{50} = 0.18$,
$\frac{8}{50} = 0.16$, $\frac{10}{50} = 0.20$, $\frac{6}{50} = 0.12$,
$\frac{10}{50} = 0.20$ (berjumlah $1$). \omterm{def:g10:stats:mean}{Rata-rata} hasilnya:
\[
\bar x = \frac{7 \times 1 + 9 \times 2 + 8 \times 3 + 10 \times 4
+ 6 \times 5 + 10 \times 6}{50}
= \frac{7 + 18 + 24 + 40 + 30 + 60}{50}
= \frac{179}{50} = 3.58 .
\]
\end{solution}

\begin{solution}{exo:g10:stats:3}
$N = 12$ (genap): \omterm{def:g10:stats:median}{mediannya} adalah \omterm{prop:g10:coordgeom:midpoint}{titik tengah} nilai ke-$6$ dan ke-$7$,
keduanya sama dengan $6$: jadi \omterm{def:g10:stats:median}{mediannya} $6$.

$Q_1$: $\frac{12}{4} = 3$, jadi $Q_1$ adalah nilai ke-$3$: yaitu $4$.

$Q_3$: $\frac{3 \times 12}{4} = 9$, jadi $Q_3$ adalah nilai ke-$9$: yaitu $8$.

Diagram kotaknya: kotak dari $4$ sampai $8$ dengan garis \omterm{def:g10:stats:median}{median} di $6$, dan
sungut dari $1$ sampai $11$.
\end{solution}

\begin{solution}{exo:g10:stats:4}
Ke-$24$ nilai itu berjumlah $24 \times 11.5 = 276$. Dengan nilai
yang baru, jumlahnya menjadi $276 + 19 = 295$ untuk $25$ siswa: \omterm{def:g10:stats:mean}{rata-rata}
barunya $\frac{295}{25} = 11.8$.
\end{solution}

\begin{solution}{exo:g10:stats:5}
Jumlah sekarang: $5 \times 12 = 60$. \omterm{def:g10:stats:mean}{Rata-rata} $13$ atas $6$ ujian menuntut
jumlah $78$, sehingga nilai ke-$6$ haruslah $78 - 60 = 18$. \omterm{def:g10:stats:mean}{Rata-rata}
$14$ menuntut $84 - 60 = 24 > 20$: jadi tidak terjangkau.
\end{solution}

\begin{solution}{exo:g10:stats:6}
\omterm{def:g10:stats:mean}{Rata-rata} lebih besar daripada $20$: misalnya
$1,\ 2,\ 3,\ 10,\ 40,\ 50,\ 60$ --- \omterm{def:g10:stats:median}{mediannya} $10$ (nilai terurut ke-$4$),
\omterm{def:g10:stats:mean}{rata-ratanya} $\frac{166}{7} \approx 23.7$.

\omterm{def:g10:stats:mean}{Rata-rata} lebih kecil daripada $6$: misalnya
$0,\ 0,\ 0,\ 10,\ 10,\ 10,\ 10$ --- \omterm{def:g10:stats:median}{mediannya} $10$, \omterm{def:g10:stats:mean}{rata-ratanya}
$\frac{40}{7} \approx 5.7$.

\omterm{def:g10:stats:mean}{Rata-ratanya} tidak dapat turun di bawah $\frac{40}{7}$: agar \omterm{def:g10:stats:median}{mediannya} (nilai
terurut ke-$4$) sama dengan $10$, nilai pada urutan $4$ sampai $7$ semuanya
harus $\geq 10$, sehingga jumlahnya paling sedikit $40$ dan \omterm{def:g10:stats:mean}{rata-ratanya}
paling sedikit $\frac{40}{7}$ --- contoh kedua itulah yang ekstrem.

Contoh ini memperlihatkan bahwa \omterm{def:g10:stats:mean}{rata-rata} dan \omterm{def:g10:stats:median}{median} sebagian besar saling
bebas: mengetahui yang satu hanya sedikit menceritakan yang lain, dan itulah
sebabnya ringkasan yang baik melaporkan keduanya.
\end{solution}

\begin{solution}{exo:g10:stats:7}
\emph{1.} \omterm{def:g10:stats:mean}{Rata-ratanya}:
$\frac{9 \times 1800 + 12000}{10} = \frac{16200 + 12000}{10} = 2820$.
\omterm{def:g10:stats:median}{Median}: setelah diurutkan, $10$ gaji itu adalah sembilan nilai $1800$ lalu
$12000$; \omterm{def:g10:stats:median}{mediannya} adalah \omterm{prop:g10:coordgeom:midpoint}{titik tengah} nilai ke-$5$ dan ke-$6$, keduanya
$1800$: jadi \omterm{def:g10:stats:median}{mediannya} $1800$.

\emph{2.} \omterm{def:g10:stats:median}{Median} ($1800$) memerikan gaji yang khas: itulah yang
sungguh-sungguh diterima $9$ dari $10$ karyawan. \omterm{def:g10:stats:mean}{Rata-rata} ($2820$) tertarik
naik oleh satu gaji tinggi dan tidak cocok dengan slip gaji siapa pun.
\end{solution}

\begin{solution}{exo:g10:stats:8}
Jumlah tingginya: $12 \times 168 + 18 \times 163 = 2016 + 2934 = 4950$ cm
untuk $30$ siswa, jadi \omterm{def:g10:stats:mean}{rata-ratanya} $\frac{4950}{30} = 165$ cm. Nilainya
lebih dekat ke \omterm{def:g10:stats:mean}{rata-rata} murid perempuan karena jumlah mereka lebih banyak
(\omterm{def:g10:stats:mean}{rata-rata} berbobot, dengan bobot $12$ dan $18$).
\end{solution}

\begin{solution}{exo:g10:stats:9}
\emph{1.} \omterm{def:g10:stats:mean}{Rata-rata} barunya adalah
\[
\bar y = \frac{y_1 + \dots + y_N}{N}
= \frac{(a x_1 + b) + \dots + (a x_N + b)}{N}
= \frac{a(x_1 + \dots + x_N) + Nb}{N}
= a \bar x + b .
\]

\emph{2.} Dengan $a = 1.8$ dan $b = 32$: \omterm{def:g10:stats:mean}{rata-ratanya}
$1.8 \times 20 + 32 = 68$ derajat Fahrenheit.
\end{solution}

\begin{solution}{pb:g10:stats:1}
\textbf{1.} \omterm{def:g10:stats:mean}{Rata-ratanya}: $\frac{2\,550}{7} \approx 364$ ribu
euro; \omterm{def:g10:stats:median}{mediannya}: harga keempat dari tujuh harga terurut, yaitu $230$.
Medianlah yang memerikan jalan itu --- enam dari tujuh rumahnya
berharga jauh dari $364$.

\textbf{2.} Tanpa rumah besar itu: \omterm{def:g10:stats:mean}{rata-ratanya} $\frac{1\,350}{6} = 225$,
\omterm{def:g10:stats:median}{mediannya} $\frac{220 + 230}{2} = 225$. \omterm{def:g10:stats:mean}{Rata-ratanya} turun sekitar
$139$; \omterm{def:g10:stats:median}{mediannya} turun $5$. Satu nilai ekstrem dapat menyeret \omterm{def:g10:stats:mean}{rata-rata}
ke mana saja; \omterm{def:g10:stats:median}{median} nyaris tidak menyadarinya --- \omterm{def:g10:stats:median}{median} bersifat \emph{kekar}.

\textbf{3.} \omterm{def:g10:stats:mean}{Rata-ratanya} $= \frac{0 \times 20 + 1 \times 25 + 2 \times
30 + 3 \times 15 + 4 \times 10}{100} = \frac{170}{100} = 1.7$.
Tidak ada keluarga yang punya $1.7$ anak, tetapi $1.7 \times 100 = 170$
\emph{memang} jumlah anak yang tepat: \omterm{def:g10:stats:mean}{rata-rata} menyandikan
jumlah keseluruhan, yang dibagi rata --- itulah bakat sejatinya.

\textbf{4.} $N = 15$: \omterm{def:g10:stats:median}{mediannya} $=$ nilai ke-8 $= 11$;
$Q_1 =$ nilai ke-4 $= 8$; $Q_3 =$ nilai ke-12 $= 14$.

\textbf{5.} \omterm{def:g10:stats:quartiles}{Jangkauannya} $19 - 4 = 15$: sebaran penuh, termasuk nilai
ekstremnya. \omterm{def:g10:stats:quartiles}{Jangkauan antarkuartilnya} $14 - 8 = 6$: lebar
paruh tengahnya --- sebaran tanpa derau nilai ekstremnya.

\textbf{6.} Pergeseran $+b$ menggerakkan \emph{setiap} nilai, sehingga
selisih antarnilai --- \omterm{def:g10:stats:quartiles}{jangkauan} dan \omterm{def:g10:stats:quartiles}{jangkauan antarkuartil} --- tidak berubah:
$(x_i + b) - (x_j + b) = x_i - x_j$. Penskalaan dengan $a$
mengalikan semua selisihnya dengan $a$: \omterm{def:g10:stats:quartiles}{jangkauan} dan \omterm{def:g10:stats:quartiles}{jangkauan antarkuartil}
terskalakan dengan $\abs a$.

\textbf{7.} Penyesuaian penjumlahan: $6 \to 9$ dan $18 \to 21$ --- siswa
lemah beruntung $50\,\%$, yang kuat $17\,\%$, dan $21$
hanya sedikit melampaui skala $20$ angka. Penyesuaian perkalian:
$6 \to 8$ dan $18 \to 24$ --- perbandingannya terpelihara, tetapi $24$
mustahil pada skala $20$: itulah masalah batas atasnya. Penyesuaian
penjumlahan menyanjung yang di bawah; penyesuaian perkalian meledakkan
yang di atas.

\textbf{8.} Jumlah nilainya: $20 \times 12 + 30 \times 10 = 540$,
sehingga \omterm{def:g10:stats:mean}{rata-rata} gabungannya $\frac{540}{50} = 10.8$: tiga puluh siswa
kelas B mengalahkan dua puluh siswa kelas A --- \omterm{def:g10:stats:mean}{rata-rata} itu bergabung
lewat \omterm{def:g10:stats:mean}{rata-rata} \emph{berbobot}, tidak pernah lewat \omterm{def:g10:stats:mean}{rata-rata} biasa.

\textbf{9.} \omterm{def:g10:stats:median}{Median} $A$: $2$; \omterm{def:g10:stats:median}{median} $B$: $5$. Daftar gabungannya
$\{0, 1, 2, 3, 10\}$: \omterm{def:g10:stats:median}{mediannya} $2$. \omterm{def:g10:stats:mean}{Rata-rata} kedua \omterm{def:g10:stats:median}{median} yang
dibobot ukurannya adalah $\frac{3 \times 2 + 2 \times 5}{5} = 3.2
\neq 2$: \omterm{def:g10:stats:median}{median} tidak mengusung jumlah keseluruhan, sehingga tidak ada rumus
penggabungan --- seluruh daftarnya harus diurutkan ulang.

\textbf{10.} \omterm{def:g10:stats:mean}{Rata-ratanya}: $\frac{31.3}{5} = 6.26$. \omterm{def:g10:stats:median}{Mediannya}: $5.4$.
\omterm{def:g10:stats:mean}{Rata-rata} terpangkasnya: $\frac{5.3 + 5.4 + 5.5}{3} = 5.4$. Satu
juri yang bersemangat (atau korup) menggeser \omterm{def:g10:stats:mean}{rata-rata} hampir satu
angka; \omterm{def:g10:stats:median}{median} dan \omterm{def:g10:stats:mean}{rata-rata} terpangkas tidak bergeming --- dari sanalah
aturan pemangkasan pada olahraga berjuri.

\textbf{11.} Pabrik A: paruh tengahnya berada dalam $0.2$ mm; pabrik B:
dalam $4$ mm. \omterm{def:g10:stats:mean}{Rata-ratanya} sama, keterandalannya berlawanan: belilah dari A.
\omterm{def:g10:stats:quartiles}{Jangkauan} antarkuartillah yang memutuskan --- \omterm{def:g10:stats:mean}{rata-rata} saja hanya
membandingkan pusat, tidak pernah keajekan.

\textbf{12.} \omterm{def:g10:stats:median}{Median} jauh di bawah \omterm{def:g10:stats:mean}{rata-rata} membocorkan adanya \emph{ekor
kanan yang panjang}: kebanyakan waktu tunggunya singkat, beberapa
membencanakan. Contohnya:
$\{5, 5, 10, 10, 95\}$ --- \omterm{def:g10:stats:median}{mediannya} $10$, \omterm{def:g10:stats:mean}{rata-ratanya}
$\frac{125}{5} = 25$.

\textbf{13.} Artinya $75\,\%$ bayi seusia itu lebih
pendek (dan $25\,\%$ lebih tinggi): itu peringkat, bukan ukuran. ``Setiap
anak di atas \omterm{def:g10:stats:mean}{rata-rata}'' mustahil --- kalau semua nilainya melampaui
\omterm{def:g10:stats:mean}{rata-ratanya}, \omterm{def:g10:stats:mean}{rata-rata} nilai itu akan melampaui dirinya sendiri:
kontradiksi. Tetapi \emph{hampir} semuanya bisa: satu anak yang sangat
pendek dapat menahan \omterm{def:g10:stats:mean}{rata-rata} di bawah semua yang lain. Di atas
\emph{\omterm{def:g10:stats:median}{median}}: tidak pernah lebih dari separuh, menurut definisinya.
Peringkat tidak dapat disanjung; \omterm{def:g10:stats:mean}{rata-rata} bisa.

\textbf{14.} Nilai yang khas: sama, \omterm{def:g10:stats:median}{mediannya} $11$ untuk keduanya.
Kehomogenan bagian tengahnya: B lebih rapat (\omterm{def:g10:stats:quartiles}{jangkauan antarkuartil} $3$
lawan $5$). Nilai ekstremnya: B lebih liar (\omterm{def:g10:stats:quartiles}{jangkauan} $17$ lawan $13$).
B adalah kelas berisi siswa yang serupa ditambah beberapa pencilan; A
tersebar merata.

\textbf{15.} (1) Apakah $5\,000$ itu \omterm{def:g10:stats:mean}{rata-rata} atau \omterm{def:g10:stats:median}{median} --- dan bolehkah
saya memperoleh \omterm{def:g10:stats:median}{mediannya}? (2) Berapa sebarannya (\omterm{def:g10:stats:quartiles}{kuartilnya}), agar saya
tahu seberapa khas ``yang khas'' itu? (3) Siapa saja yang terhitung ---
ukuran sampel dan cara pemilihannya (pelajaran jajak pendapat pada soal
akhir pekan tentang cara data berdusta di jilid sebelumnya)?

\textbf{16.} Misalnya $100, 150, 200, 250, 800$: \omterm{def:g10:stats:median}{mediannya}
$200$, \omterm{def:g10:stats:mean}{rata-ratanya} $\frac{1\,500}{5} = 300$.

\textbf{17.} $\sum (x_i - \bar x) = \sum x_i - N \bar x =
N\bar x - N\bar x = 0$. Periksa: $-200 - 150 - 100 - 50 + 500 =
0$. \omterm{def:g10:stats:mean}{Rata-rata} adalah titik tempat simpangannya seimbang --- yaitu
titik berat datanya (titik $G$ pada soal akhir pekan tentang titik berat
di jilid sebelumnya, dalam satu dimensi).

\textbf{18.} Di \omterm{def:g10:stats:median}{median} $4$: $3 + 1 + 0 + 5 + 9 = 18$ km. Di
\omterm{def:g10:stats:mean}{rata-rata} $6$: $5 + 3 + 2 + 3 + 7 = 20$ km. Di $5$:
$4 + 2 + 1 + 4 + 8 = 19$ km. Medianlah yang menang --- bergerak menjauh
darinya selalu melewati lebih banyak toko daripada yang didekatinya,
sehingga jarak totalnya hanya dapat membesar.

\textbf{19.} $0.3^{10} \approx 6 \times 10^{-6}$: sekitar enam
penerbang \emph{per sejuta}. Pada $4\,000$ penerbang, cacah
harapannya $0.02$ --- yaitu nol, seperti yang ditemukan letnan Daniels.
Angkatan Udara berhenti merancang untuk lelaki \omterm{def:g10:stats:mean}{rata-rata} dan membuat
segalanya dapat disetel: kursi, pedal, sabuk --- merancang untuk
\emph{sebarannya}, bukan untuk pusatnya.

\textbf{20.} Untuk jumlah keseluruhan: \emph{\omterm{def:g10:stats:mean}{rata-rata}}. Untuk individu yang
khas: \emph{\omterm{def:g10:stats:median}{median}}. Untuk keadilan sebarannya: \emph{\omterm{def:g10:stats:quartiles}{kuartil}} (dan
\omterm{def:g10:stats:quartiles}{jangkauan antarkuartil}). Untuk pencemaran data: \emph{\omterm{def:g10:stats:mean}{rata-rata} terpangkas}.
Lalu baris penutupnya: orang \omterm{def:g10:stats:mean}{rata-rata} itu tidak ada --- tetapi orang
\omterm{def:g10:stats:median}{median} hampir ada.

\end{solution}
